\chapter{المجموعات والبنى}\label{ch:b2:structures}

يشحذ هذا الفصل الافتتاحي الأسس الموضوعة في مجلد السنة الأولى ويحوّلها
إلى أدوات عمل يومية: حساب المجموعات ومجموعات القسمة، ومقارنة المجموعات
غير المنتهية (قابلية العد، كانتور--برنشتاين)، والنظرية البنيوية للزمر
والحلقات --- مبرهنة لاغرانج، والزمرة المتناظرة مع إشارتها، والمثاليات
ومبرهنة البواقي الصينية. كل ما هنا يُستعمل بلا انقطاع في بقية الكتاب:
الإشارة تبني المحدد
(\cref{ch:b2:linalg})، \omterm{def:b2:structures:quotientring}{وحلقات القسمة} تُدير الحساب، وقابلية
العد تقوم تحت الطوبولوجيا والاحتمال معًا.

\section{المجموعات والتطبيقات ومجموعات القسمة}

نستعمل بحرية لغة المجموعات والتطبيقات وعلاقات التكافؤ والترتيب التي
أُرسيت في مجلد السنة الأولى. وهناك تحسينان يستحقان صياغة صريحة.

\begin{proposition}[صور العائلات وصورها العكسية]\label{prop:b2:structures:images}
ليكن $f \colon E \to F$ ولتكن $(A_i)_{i \in I}$ و$(B_j)_{j \in J}$
عائلتين من أجزاء $E$ و$F$ على التوالي. عندئذٍ
\[
f^{-1}\Bigl(\bigcup_j B_j\Bigr) = \bigcup_j f^{-1}(B_j),
\qquad
f^{-1}\Bigl(\bigcap_j B_j\Bigr) = \bigcap_j f^{-1}(B_j),
\qquad
f^{-1}(F \setminus B) = E \setminus f^{-1}(B),
\]
\[
f\Bigl(\bigcup_i A_i\Bigr) = \bigcup_i f(A_i),
\qquad
f\Bigl(\bigcap_i A_i\Bigr) \subseteq \bigcap_i f(A_i)
\quad (\text{مع المساواة إذا كان } f).
\] متباينًا
\end{proposition}

\begin{proof}
كل متطابقة هنا ليست إلا فكًّا للتعاريف؛ فمثلًا $x \in
f^{-1}(\bigcap B_j) \iff f(x) \in B_j$ لكل $j$ $\iff x \in
f^{-1}(B_j)$ لكل $j$. أما متطابقات الصورة وانعدام
المساواة في حالة التقاطع (مع تصحيحها بالتباين) فقد بُرهنت في مجلد
السنة الأولى من أجل مجموعتين؛ والحجج نفسها تصلح للعائلات.
\end{proof}

\begin{example}[حيث يكون احتواء الصور تامًا]\label{ex:b2:structures:imagestrict}
لنأخذ $f \colon \R \to \R$، $f(x) = x^2$، مع $A_1 =
\intcc{-1}{0}$ و$A_2 = \intcc{0}{1}$. عندئذٍ
\[
f(A_1 \cap A_2) = f(\{0\}) = \{0\},
\qquad
f(A_1) \cap f(A_2) = \intcc{0}{1} \cap \intcc{0}{1} =
\intcc{0}{1} :
\]
فالاحتواء في \cref{prop:b2:structures:images} تامّ إلى أقصى
حدّ --- إذ إن نقطتي الصورة العكسية $\pm x$ لقيمة مشتركة تقعان في
$A_i$ مختلفين. والتباين هو بالضبط ما يمنع هذا
الانفصال، ولهذا تُحقق الصور العكسية (التي لا تدمج نقطتين أبدًا)
المتطابقات الأربع دون شرط، بينما تفقد الصور تلك المتعلقة
بالتقاطع. قاعدة عملية للكتاب كله: مرّر \emph{الصور العكسية} عبر
عمليات المجموعات بلا تردد، وعامِل الصور بحذر.
\end{example}

\begin{definition}[مجموعة القسمة]\label{def:b2:structures:quotient}
لتكن $\mathcal{R}$ علاقة تكافؤ على $E$. تُسمى
\emph{مجموعة القسمة}\index{مجموعة القسمة} $E/\mathcal{R}$ مجموعة
أصناف التكافؤ؛ والتطبيق الشامل $\pi \colon E \to E/\mathcal{R}$،
$x \mapsto \mathrm{cl}(x)$، هو \emph{الإسقاط
القانوني}\index{الإسقاط القانوني}.

\emph{الخاصية الشمولية (التحليل إلى عوامل):} إذا كان $f \colon E \to F$
\emph{متوافقًا} مع $\mathcal{R}$ (أي $x \mathbin{\mathcal{R}}
y \implies f(x) = f(y)$)، فهناك تطبيق وحيد $\overline f
\colon E/\mathcal{R} \to F$ يحقق $f = \overline f \circ \pi$.
\end{definition}

\begin{proof}[برهان الخاصية الشمولية]
الوحدانية: الشرط $f = \overline f \circ \pi$ يُكتب
\[
\overline f\bigl(\mathrm{cl}(x)\bigr) = f(x)
\qquad (x \in E),
\]
وبما أن $\pi$ شامل، فإن كل عنصر من
$E/\mathcal{R}$ هو $\mathrm{cl}(x)$ ما: فقيم
$\overline f$ كلها مفروضة. الوجود: نأخذ الصيغة المعروضة
\emph{تعريفًا} للتطبيق $\overline f$؛ وهي غير ملتبسة
بالضبط بفضل التوافق --- إذ إن $\mathrm{cl}(x) =
\mathrm{cl}(y)$ يعطي $x \mathbin{\mathcal{R}} y$، ومنه $f(x) =
f(y)$ فتتفق القيمتان المرشحتان --- وهي تُحلّل
$f$ بحكم الإنشاء. ولاحظ تقسيم العمل: شمول
$\pi$ يعطي الوحدانية، والتوافق يعطي الوجود.
\end{proof}

\begin{example}
$\Z/n\Z$ هي قسمة $\Z$ على التوافق ترديد $n$؛ وعمليات
التحقق من حسن التعريف في مجلد السنة الأولى كانت حالات خاصة من
الخاصية الشمولية. فمجموعات القسمة تحوّل ``الإنشاءات المتوافقة على
الممثلين'' إلى تطبيقات صريحة --- وهذا ما نستعمله باستمرار فيما يلي.
\end{example}

\section{قابلية العد والقوة}

\begin{definition}[تساوي القوة، قابلية العد]\label{def:b2:structures:countable}
تكون مجموعتان \emph{متساويتي القوة}\index{تساوي القوة} إذا وُجد تقابل
بينهما. وتكون مجموعة \emph{قابلة للعد}\index{مجموعة قابلة للعد} إذا
كانت مساوية القوة للمجموعة $\N$ (يُدرج بعض المؤلفين المجموعات المنتهية؛ نقول
نحن \emph{قابلة للعد على الأكثر} بمعنى ``منتهية أو قابلة للعد'').
\end{definition}

\begin{proposition}[خواص الاستقرار]\label{prop:b2:structures:countablestable}
\begin{enumerate}
  \item كل جزء غير منته من $\N$ قابل للعد؛ ومجموعة تكون \omterm{def:b2:structures:countable}{قابلة للعد}
        على الأكثر إذا وفقط إذا انغرست في $\N$ إذا وفقط إذا كانت
        خالية أو صورة شاملة للمجموعة $\N$.
  \item $\N \times \N$ \omterm{def:b2:structures:countable}{قابلة للعد}؛ وجداء مجموعتين قابلتين للعد على
        الأكثر قابل للعد على الأكثر.
  \item اتحاد قابل للعد على الأكثر من مجموعات \omterm{def:b2:structures:countable}{قابلة للعد} على الأكثر
        هو قابل للعد على الأكثر.
  \item $\Z$ و$\Q$ قابلتان للعد.
\end{enumerate}
\end{proposition}

\begin{proof}
(1) نُرتّب جزءًا غير منته $A \subseteq \N$ بأخذ الأصغر تكرارًا: $a_0 =
\min A$، $a_{k+1} = \min\,(A \setminus \{a_0, \dots, a_k\})$
(غير خالية لأن $A$ غير منتهية)؛ فالتطبيق $k \mapsto a_k$ متزايد تمامًا،
ومتباين، وشامل على $A$ (فكل $a \in A$
يفوق عددًا منتهيًا فقط من عناصر $A$، ومن ثم يُبلَغ). وإذا انغرست $E$
في $\N$ بواسطة $\varphi$، فإن $E$ مساوية القوة للمجموعة $\varphi(E) \subseteq \N$: أي
منتهية أو \omterm{def:b2:structures:countable}{قابلة للعد}. وإذا كان $s \colon \N \to
E$ شاملًا، فإن $x \mapsto \min s^{-1}(\{x\})$ يغرس $E$
في $\N$.

(2) التطبيق $(p, q) \mapsto 2^p(2q + 1) - 1$ تقابل $\N^2
\to \N$ (فلكل عدد صحيح موجب تفكيك وحيد إلى
فردي وقوة للعدد 2 على الصورة $2^p m$ مع $m$ فردي، بحكم وحدانية
التفكيك). أما الجداءات: فنركّب الغرسات.

(3) لتكن مجموعات $E_n$ مع تطبيقات شاملة $s_n \colon \N \to E_n$
(ولا ضير إذا كانت بعض $E_n$ منتهية: نكرر القيم)، فالتطبيق $(n, k)
\mapsto s_n(k)$ شامل من \omterm{def:b2:structures:countable}{المجموعة القابلة للعد} $\N^2$ على
$\bigcup E_n$.

(4) $\Z = \N \cup (-\N^*)$: اتحاد قابل للعد. و$\Q$ صورة شاملة
للمجموعة $\Z \times \N^*$ (تطبيق الكسر)، ومن ثم \omterm{def:b2:structures:countable}{قابلة للعد} على
الأكثر، وهي غير منتهية.
\end{proof}

\begin{example}[دالة ازدواج، محسوبة]\label{ex:b2:structures:pairing}
التقابل $(p, q) \mapsto 2^p(2q + 1) - 1$ الوارد في البرهان
جدير بأن نراه في العمل. وهذه قيمه الأولى:
\[
\begin{array}{c|ccccc}
 & q = 0 & q = 1 & q = 2 & q = 3 & q = 4\\
\hline
p = 0 & 0 & 2 & 4 & 6 & 8\\
p = 1 & 1 & 5 & 9 & 13 & 17\\
p = 2 & 3 & 11 & 19 & 27 & 35\\
p = 3 & 7 & 23 & 39 & 55 & 71
\end{array}
\]
يجمع السطر $p$ الأعداد الصحيحة $n$ التي يقبل من أجلها $n + 1$ القسمة
تمامًا على $2^p$: فكل عدد طبيعي يظهر مرة واحدة بالضبط.
وفكّ الترميز صريح كترميزه: من أجل $n = 43$، نفكك $n + 1
= 44 = 2^2\cdot 11 = 2^2(2\cdot5 + 1)$، ومنه $(p, q) = (2, 5)$.
والفكرة الختامية: براهين قابلية العد كثيرًا ما تكون
\emph{خوارزميات} متخفية --- وهي هنا ``أخرِج قوى العدد 2''.
\end{example}

\begin{example}[الأعداد الجبرية قابلة للعد]\label{ex:b2:structures:algebraic}
يكون عدد عقدي \emph{جبريًا} إذا أعدم كثير حدود غير معدوم بمعاملات
ناطقة. ومجموعة الأعداد الجبرية $\overline\Q$ \omterm{def:b2:structures:countable}{قابلة للعد}: فكثيرات الحدود
من الدرجة $\leq d$ على $\Q$ تنغرس في $\Q^{d+1}$، وهو جداء منته
لمجموعات \omterm{def:b2:structures:countable}{قابلة للعد}
(\cref{prop:b2:structures:countablestable}~(2))؛ والاتحاد على
$d$ يُعدّد كثيرات الحدود الناطقة غير المعدومة على الصورة $P_0, P_1,
P_2, \dots$؛ ولكل $P_k$ عدد منته من الجذور؛ و
\[
\overline\Q = \bigcup_{k \in \N}\ \{\text{جذور } P_k\}
\]
هو اتحاد قابل للعد لمجموعات منتهية
(\cref{prop:b2:structures:countablestable}~(3))، وهو غير منته لأنه
يحتوي $\Q$. وباقتران ذلك بعدم قابلية $\R$ للعد
(\cref{thm:b2:structures:cantor} أدناه)، نبرهن --- دون
إبراز عدد واحد --- على وجود الأعداد المتسامية وعلى أنها تشكّل
الأغلبية غير \omterm{def:b2:structures:countable}{القابلة للعد}: إنها حجة العدّ عند كانتور سنة
1874، وجودٌ بالقوة وحدها.
\end{example}

\begin{theorem}[كانتور؛ عدم قابلية $\R$ للعد]\label{thm:b2:structures:cantor}
\begin{enumerate}
  \item من أجل كل مجموعة $E$، لا يوجد تطبيق شامل $E \to
        \mathcal{P}(E)$.
  \item $\R$ \emph{غير} \omterm{def:b2:structures:countable}{قابلة للعد}.
\end{enumerate}
\end{theorem}

\begin{proof}
(1) بُرهن عليه في مجلد السنة الأولى (المجموعة القطرية $D = \{x : x
\notin f(x)\}$).

(2) لنفترض أن $(x_n)_{n \in \N}$ يُعدّد $\R$. نبني قطعًا متداخلة
$I_0 \supseteq I_1 \supseteq \dots$ بحيث $\abs{I_n} = 3^{-n}$ و
$x_n \notin I_n$: نقسم القطعة الجارية إلى ثلاثة أثلاث مغلقة؛
فأحد الأثلاث على الأقل يتجنب $x_n$ (لأن نقطة تلتقي اثنين على الأكثر
من الثلاثة). ومبرهنة القطع المتداخلة (بأطراف متجاورة) تعطي
$\ell \in \bigcap_n I_n$؛ لكن $\ell = x_N$ من أجل $N$ ما، و$x_N
\notin I_N$: تناقض.
\end{proof}

\begin{theorem}[كانتور--برنشتاين]\label{thm:b2:structures:cantorbernstein}
إذا انغرست $E$ في $F$ وانغرست $F$ في $E$، فإن $E$ و$F$
متساويتا القوة.\index{مبرهنة كانتور--برنشتاين}
\end{theorem}

\begin{proof}
ليكن $f \colon E \to F$ و$g \colon F \to E$ غرستين. من أجل كل نقطة (من $E$ أو $F$)،
نتتبع \emph{سلسلة أسلافها} من الصور العكسية المتتابعة،
$x \mapsto g^{-1}(x) \mapsto
f^{-1}(g^{-1}(x)) \mapsto \dots$ --- وكل خطوة معرَّفة ما دامت النقطة الجارية تنتمي إلى صورة
الغرسة المعنية، وهي عندئذٍ وحيدة بحكم التباين. وهناك ثلاثة مصائر
يستبعد بعضها بعضًا: تتوقف السلسلة عند نقطة من $E \setminus g(F)$ (\emph{المنشأ في
$E$})، أو تتوقف عند نقطة من $F \setminus f(E)$ (\emph{المنشأ في
$F$})، أو لا تتوقف أبدًا. وهذا يقسم $E = E_E \cup E_F \cup
E_\infty$ و$F = F_E \cup F_F \cup F_\infty$ حسب
المنشأ.

ولنلاحظ الآن: $f$ يرسل $E_E$ \emph{على} $F_E$ --- فسلسلة
$f(x)$ هي سلسلة $x$ مسبوقة بخطوة واحدة، ومن ثم يتطابق المنشآن؛
ولكل $y \in F_E$ سلسلة فيها خطوة واحدة على الأقل (لأن منشأه
يقع في $E$)، ومنه $y = f(x)$ مع $x \in E_E$. والحجة نفسها
تعطي تقابلين $f \colon E_\infty \to F_\infty$ و$g \colon F_F
\to E_F$. وباللصق، فإن
\[
h(x) =
\begin{cases}
f(x) & \text{إذا كان } x \in E_E \cup E_\infty,\\
g^{-1}(x) & \text{إذا كان } x \in E_F,
\end{cases}
\]
تقابل من $E$ على $F = F_E \cup F_\infty \cup F_F$: فهو تقابل على كل قطعة،
والقطع الثلاث في المجموعة الهدف منفصلة مثنى مثنى.
\end{proof}

\begin{example}\label{ex:b2:structures:cbexample}
$\intoo{0}{1}$ و$\intcc{0}{1}$ متساويتا القوة: فالمطابق
يغرس في اتجاه، و$x \mapsto \frac{x + 1}{3}$ في الاتجاه الآخر؛ والمبرهنة
تصنع التقابل (غير المتصل بالضرورة). وكذلك
$\R$ و$\intoo{0}{1}$ (بتقابلات من نمط $\tanh$) و
$\mathcal{P}(\N)$ (بالنشر الثنائي، \cref{exo:b2:structures:3})
كلها متساوية القوة: إنها ``قوة المتصل''.
\end{example}

\begin{example}[القطعة والمربع]\label{ex:b2:structures:squareseg}
$\intcc{0}{1}$ و$\intcc{0}{1}^2$ متساويتا القوة --- فالبعد غير مرئي
للقوة. إحدى الغرستين بديهية: $x
\mapsto (x, 0)$. أما الأخرى فنرسل $(x, y)$
إلى العدد الحقيقي الذي تتشابك أرقامه العشرية مع أرقام $x$ و$y$،
\[
(0.x_1x_2x_3\dots,\ 0.y_1y_2y_3\dots)
\;\longmapsto\; 0.x_1y_1x_2y_2x_3y_3\dots,
\]
مع اختيار النشر الذي لا ينتهي بالرقم $9$ المتكرر من أجل كل إحداثي:
وبهذا الاصطلاح تحدد أرقام الصورة أرقام $x$ و$y$، فالتطبيق
متباين (ولا يلزم أن يكون شاملًا --- إذ لا تحتوي الصور مثلًا على
أرقام في المواضع الفردية تنتهي إلى $9$ --- ولا ضير في ذلك).
ومبرهنة كانتور--برنشتاين (\cref{thm:b2:structures:cantorbernstein})
تُركّب تقابلًا حقيقيًا. أما الاتصال فميؤوس منه بالطبع: لا يمكن أن
يوجد تقابل متصل بينهما --- وفصول الفضاءات المترية تشرح السبب
(الترابط يميز المستقيم عن المستوي، \cref{ch:b2:metric}).
\end{example}

\section{الزمر}

\begin{definition}[الزمرة الجزئية المولَّدة؛ الرتبة]\label{def:b2:structures:generated}
لتكن $G$ زمرة وليكن $A \subseteq G$. الزمرة الجزئية \emph{المولَّدة}
بالجزء $A$، ويُرمز إليها $\langle A \rangle$، هي أصغر زمرة جزئية
تحتوي $A$ --- وهي بالتحديد مجموعة كل الجداءات المنتهية لعناصر $A$
ومقلوباتها. وتكون زمرة \emph{دورية}\index{زمرة دورية} إذا كانت
مولَّدة بعنصر واحد: $\langle a\rangle = \{a^k : k \in \Z\}$.
و\emph{رتبة}\index{رتبة!رتبة عنصر} العنصر $a \in G$ هي
$\operatorname{ord}(a) = \abs{\langle a \rangle}$ (وقد تكون
غير منتهية)؛ وإذا كانت منتهية فهي أصغر $n \geq 1$ يحقق $a^n = e$،
ويكون $a^k = e \iff \operatorname{ord}(a) \mid k$.
\end{definition}

\begin{proof}[برهان توصيف الرتبة]
إذا وُجد $a^m = e$ يحقق $m \geq 1$، فليكن $n \geq 1$ أصغر عدد يحقق $a^n
= e$. عندئذٍ تكون العناصر $e, a, \dots, a^{n-1}$ مختلفة مثنى مثنى
(إذ إن $a^{i} = a^{j}$ مع $0 \leq i < j < n$ يعطي $a^{j-i} = e$،
وهذا ينقض الأصغرية)، وكل $a^k$ يؤول إلى أحدها بالقسمة الإقليدية
$k = nq + r$: فعدد عناصر $\langle a\rangle$ هو $n$ بالضبط، و$a^k = a^r = e \iff r = 0 \iff n \mid k$. وإذا لم تكن أي قوة
محايدة، فإن جميع $a^k$ ($k \in \Z$) مختلفة (بحجة القسمة نفسها)
والرتبة غير منتهية.
\end{proof}

\begin{theorem}[لاغرانج]\label{thm:b2:structures:lagrange}
لتكن $G$ زمرة منتهية ولتكن $H$ زمرة جزئية منها. عندئذٍ يقسم $\abs H$
العدد $\abs G$.\index{مبرهنة لاغرانج} وعلى الخصوص تقسم رتبة كل
عنصر $\abs G$، ويكون $a^{\abs G} = e$ لكل $a \in G$.
\end{theorem}

\begin{proof}
العلاقة $x \sim y \iff x^{-1}y \in H$ علاقة تكافؤ
(انعكاسية: $e \in H$؛ تناظرية: بالمقلوبات؛ متعدية: بالجداءات).
وصنف $x$ هو \emph{الصنف الجانبي الأيسر} $xH = \{xh : h \in H\}$،
والتطبيق $h \mapsto xh$ تقابل $H \to xH$ (مقلوبه $y \mapsto
x^{-1}y$): فلجميع الأصناف $\abs H$ عنصرًا. والأصناف تقسم $G$
(مبرهنة التقسيم العامة في مجلد السنة الأولى)، ومنه $\abs G =
\abs H \times (\text{عدد الأصناف الجانبية})$. أما من أجل عنصر: فنطبق هذا
على $H = \langle a\rangle$؛ فينتج $a^{\abs G} = (a^{\operatorname{ord}
a})^{\abs G / \operatorname{ord} a} = e$.
\end{proof}

\begin{example}[الأصناف الجانبية في العمل: $A_3$ داخل $\mathfrak{S}_3$]\label{ex:b2:structures:cosets}
لنأخذ $G = \mathfrak{S}_3$ (رتبتها $6$) و$H = A_3 =
\{\mathrm{id},\ (1\,2\,3),\ (1\,3\,2)\}$. الأصناف الجانبية اليسرى هي
\[
H = \{\mathrm{id},\ (1\,2\,3),\ (1\,3\,2)\},
\qquad
(1\,2)H = \{(1\,2),\ (2\,3),\ (1\,3)\} :
\]
أي صنفان من ثلاثة عناصر يقسمان $G$، تمامًا كما يقتضي
العدّ $\abs G = \abs H \times (\text{عدد الأصناف الجانبية})$
--- وهو بوضوح التقسيم إلى تبديلات زوجية وفردية. ولاحظ أن
$(1\,3)H = (1\,2)H$ رغم أن $(1\,3) \neq
(1\,2)$: فالأصناف الجانبية \emph{أصناف} لا تُسمّى
بممثليها، و$x^{-1}y \in H$ هي المقارنة المشروعة الوحيدة. وهذه الصورة
ذات الصنفين هي الصورة العامة للإشارة: فالزمرة $A_n$ وصنفها الجانبي
الوحيد المرافق يقسمان $\mathfrak{S}_n$ نصفين، وهذا ما تعدّ به مسألة نهاية
الأسبوع الوضعيات القابلة للبلوغ في الأحجية.
\end{example}

\begin{example}\label{ex:b2:structures:lagrangeapp}
فائدتان مباشرتان. \emph{الزمر ذات الرتبة الأولية دورية:}
إذا كان $\abs G = p$ أوليًا و$a \neq e$، فإن
$\operatorname{ord}(a)$ يقسم $p$ ولا يساوي $1$، ومنه فهو $p$:
أي $\langle a\rangle = G$. \emph{شبكة الزمر الجزئية للزمرة $\Z/12\Z$:}
حسب \cref{prop:b2:structures:cyclic} أدناه، هناك زمرة جزئية
واحدة بالضبط لكل قاسم من قواسم $12$ --- رتبها $1, 2, 3, 4, 6, 12$،
مولَّدة على التوالي بالأصناف $\overline 0$ و$\overline 6$ و$\overline
4$ و$\overline 3$ و$\overline 2$ و$\overline 1$. والتنبيه
الختامي: \emph{عكس} مبرهنة لاغرانج خاطئ عمومًا --- فرتبة $A_4$
هي $12$ ولا تملك زمرة جزئية رتبتها $6$، كما نبرهن في مسألة نهاية
الأسبوع في هذا الفصل (\cref{pb:b2:structures:1}، السؤال 14).
فلاغرانج يُقيّد الرتب الممكنة، لكنه لا يَعِد بها.
\end{example}

\begin{omfigure}
\begin{tikzpicture}[scale=1.05, every node/.style={font=\small}]
  \node (G) at (0,3)  {$\Z/12\Z = \langle\overline1\rangle$};
  \node (A) at (-1.6,2) {$\langle\overline2\rangle$};
  \node (B) at (1.6,2)  {$\langle\overline3\rangle$};
  \node (C) at (-1.6,1) {$\langle\overline4\rangle$};
  \node (D) at (1.6,1)  {$\langle\overline6\rangle$};
  \node (E) at (0,0)  {$\{\overline0\}$};
  \draw[omDef] (G) -- (A); \draw[omDef] (G) -- (B);
  \draw[omDef] (A) -- (C); \draw[omDef] (A) -- (D);
  \draw[omDef] (B) -- (D); \draw[omDef] (C) -- (E);
  \draw[omDef] (D) -- (E);
  \node[omThm, right=2pt] at (0.1,3) {\ الرتبة $12$};
  \node[omThm, left=2pt]  at (-2.1,2) {الرتبة $6$\ };
  \node[omThm, right=2pt] at (2.1,2) {\ الرتبة $4$};
  \node[omThm, left=2pt]  at (-2.1,1) {الرتبة $3$\ };
  \node[omThm, right=2pt] at (2.1,1) {\ الرتبة $2$};
\end{tikzpicture}

{\small شبكة الزمر الجزئية للزمرة $\Z/12\Z$: زمرة جزئية واحدة لكل
قاسم من قواسم $12$ (\cref{prop:b2:structures:cyclic})، مع حافة
حين تحتوي إحداهما الأخرى بدليل أولي. والاحتواءات تسير
\emph{عكس} قابلية القسمة للمولّد: $\langle\overline
4\rangle \subseteq \langle\overline2\rangle$ لأن $4$
مضاعف للعدد $2$.}
\end{omfigure}

\begin{proposition}[الزمر الدورية]\label{prop:b2:structures:cyclic}
لتكن $G = \langle a \rangle$ دورية رتبتها $n$.
\begin{enumerate}
  \item $G$ متماثلة مع $(\Z/n\Z, +)$، بواسطة $\overline k \mapsto
        a^k$.
  \item كل زمرة جزئية من $G$ دورية؛ ولكل قاسم $d \mid n$
        توجد زمرة جزئية وحيدة رتبتها $d$، وهي $\langle
        a^{n/d}\rangle$.
  \item يولّد $a^k$ الزمرة $G$ إذا وفقط إذا كان $\gcd(k, n) = 1$: فعدد مولّدات
        $G$ يساوي $\varphi(n)$ (دالة أويلر\index{دالة أويلر}).
\end{enumerate}
\end{proposition}

\begin{proof}
(1) التطبيق $k \mapsto a^k$ من $\Z$ على $G$ متوافق مع
الترديد بترديد $n$ ($a^{k} = a^{k'} \iff n \mid k - k'$، حسب توصيف
الرتبة)؛ والخاصية الشمولية
(\cref{def:b2:structures:quotient}) تعطي تشاكلًا تقابليًا
معرَّفًا جيدًا انطلاقًا من $\Z/n\Z$.

(2) لتكن $H \leq G$ غير محايدة وليكن $m$ أصغر $\geq 1$ يحقق $a^m \in
H$. تُبيّن القسمة الإقليدية أن $H = \langle a^m\rangle$ (من أجل $a^k \in
H$: يفرض $k = mq + r$ أن $a^r \in H$، ومنه $r = 0$)، وأن $m \mid n$
(بقسمة $n$ على $m$: $a^{n \bmod m} \in H$). عندئذٍ $\abs H = n/m$؛
وأخذ $m = n/d$ يحقق كل قاسم $d$. أما الوحدانية: فكل زمرة جزئية
رتبتها $d$ هي، حسب ما تقدم، من الشكل $\langle a^m \rangle$
مع $n/m = d$ --- ومن ثم يُفرض $m = n/d$ وتتحدد الزمرة الجزئية.

(3) ندّعي أن $\operatorname{ord}(a^k) = \frac{n}{\gcd(k, n)}$.
لنكتب $d = \gcd(k, n)$. من أجل أي $m \geq 1$، يعطي توصيف
الرتبة في \cref{def:b2:structures:generated} سلسلة التكافؤات
\[
(a^k)^m = e
\iff n \mid km
\iff \frac{n}{d} \,\Big|\, \frac{k}{d}\,m
\iff \frac{n}{d} \,\Big|\, m ,
\]
والخطوة الأخيرة بمبرهنة غاوس، لأن $\frac nd$ و$\frac kd$
أوليان فيما بينهما. وأصغر $m$ كهذا هو $\frac nd$:
أي $\operatorname{ord}(a^k) = \frac n{\gcd(k,n)}$، وهو يساوي $n$
إذا وفقط إذا $\gcd(k, n) = 1$. وهناك $\varphi(n)$ صنفًا $k$ كهذه
بترديد $n$.
\end{proof}

\section{الزمرة المتناظرة}

\begin{definition}\label{def:b2:structures:sn}
$\mathfrak{S}_n$ هي زمرة تبديلات $\intint{1}{n}$
(رتبتها $n!$). و\emph{الدورة}\index{دورة} $(a_1\,a_2\,\cdots\,a_k)$
ترسل $a_1 \mapsto a_2 \mapsto \dots \mapsto a_k \mapsto a_1$
وتثبّت كل ما عداها؛ و$k$ هو \emph{طولها}، والدورة ذات الطول $2$ تُسمى
\emph{مبادلة}\index{مبادلة}. وتكون دورتان
\emph{منفصلتين} إذا كان حاملاهما (أي مجموعتا النقاط غير الثابتة)
منفصلين.
\end{definition}

\begin{theorem}[التفكيك إلى دورات]\label{thm:b2:structures:cycles}
كل تبديلة $\sigma \neq \mathrm{id}$ جداء دورات منفصلة مثنى مثنى، وذلك بكيفية
وحيدة إلى حدّ ترتيب العوامل. والدورات المنفصلة تتبادل، ورتبة
$\operatorname{ord}(\sigma)$ هي المضاعف المشترك الأصغر للأطوال.
\end{theorem}

\begin{proof}
لننظر في علاقة ``المدار'' على حامل $\sigma$: $x \sim y$
إذا وفقط إذا $y = \sigma^k(x)$ من أجل $k \in \Z$ ما --- وهي علاقة
تكافؤ. وكل صنف $\{x, \sigma(x), \dots, \sigma^{k-1}(x)\}$
(منته، فتعود التكرارات دوريًا --- إذ يجب أن يعود أول تكرار
إلى $x$ بحكم التباين) يحمل \omterm{def:b2:structures:sn}{الدورة} $(x\ \sigma(x)\
\cdots\ \sigma^{k-1}(x))$، و$\sigma$ هو جداء هذه
الدورات: فعلى كل مدار لا تؤثر إلا \omterm{def:b2:structures:sn}{الدورة} الموافقة.
أما الوحدانية: فكل تفكيك إلى دورات منفصلة يعيد إنتاج المدارات
بالضبط (إذ يجب أن تكون \omterm{def:b2:structures:sn}{الدورة} المارّة بالنقطة $x$ هي $(x\ \sigma(x)\ \cdots)$).
والدورات المنفصلة تتبادل لأنها تحرّك نقاطًا منفصلة؛ وينتج قول
الرتبة لأن $\sigma^m = \mathrm{id}$ إذا وفقط إذا كانت قوة
$m$ لكل دورة محايدة (بحكم الانفصال)، إذا وفقط إذا قسم كل طول $m$.
\end{proof}

\begin{example}[نمط الدورات بوصفه إحصاءً]\label{ex:b2:structures:cycletype}
كم عدد تبديلات $\mathfrak{S}_9$ ذات نمط الدورات
$(4, 3, 2)$ --- دورة من الطول $4$، ودورة من الطول $3$، \omterm{def:b2:structures:sn}{ومبادلة}؟
نختار الحوامل والترتيبات الدورية:
\[
\frac{9!}{4\cdot 3\cdot 2}
= \frac{362\,880}{24} = 15\,120 :
\]
نرصّ الرموز التسعة في سطر ($9!$ طريقة)، ونُقوّس الأربعة
الأولى ثم الثلاثة التالية ثم الأخيرين في دورات، ونقسم على
الدورانات داخل كل قوس ($4$ و$3$ و$2$ منها) التي
تعطي التبديلة نفسها. (والأطوال هنا \emph{مختلفة}،
فلا قسمة إضافية؛ أما الأطوال المتساوية فتقتضي أيضًا
القسمة على تبديلات الأقواس المتساوية.) ورتبة كل تبديلة
كهذه $\operatorname{lcm}(4,3,2) = 12$ وإشارتها
$(-1)^3(-1)^2(-1)^1 = +1$
(\cref{thm:b2:structures:cycles} ومبرهنة الإشارة
أدناه). تقسيم واحد للعدد $9$، وصنف تقارن واحد، وإحصاء واحد
--- فتوافقيات $\mathfrak{S}_n$ هي حساب التقسيمات.
\end{example}

\begin{theorem}[الإشارة]\label{thm:b2:structures:signature}
يوجد تشاكل زمر وحيد $\varepsilon \colon
\mathfrak{S}_n \to \{\pm 1\}$ (من أجل $n \geq 2$) يأخذ القيمة $-1$
على المبادلات: وهو \emph{الإشارة}\index{إشارة}. زيادة على ذلك،
$\varepsilon(\sigma) = (-1)^{I(\sigma)}$ حيث $I(\sigma)$ هو
عدد \emph{الانقلابات} (أي الأزواج $i < j$ التي $\sigma(i) >
\sigma(j)$)، وإشارة
\omterm{def:b2:structures:sn}{الدورة} ذات الطول $k$ هي $(-1)^{k-1}$، ورتبة \emph{الزمرة المتناوبة}
$A_n = \ker\varepsilon$ هي $\frac{n!}{2}$.
\end{theorem}

\begin{proof}
\emph{الوجود.} من أجل $\sigma \in \mathfrak{S}_n$ نضع
\[
\varepsilon(\sigma)
= \prod_{1 \leq i < j \leq n}
\frac{\sigma(j) - \sigma(i)}{j - i} .
\]
تتضاعف القيم المطلقة للعوامل فيعطي جداؤها $1$ (لأن الأزواج غير
المرتبة $\{\sigma(i), \sigma(j)\}$ تجري على كل الأزواج)، ومنه
$\varepsilon(\sigma) = (-1)^{I(\sigma)} \in \{\pm1\}$. أما التشاكل: فمن أجل
$\sigma, \tau$،
\[
\varepsilon(\sigma\tau)
= \prod_{i<j} \frac{\sigma(\tau(j)) - \sigma(\tau(i))}{j - i}
= \prod_{i<j} \frac{\sigma(\tau(j)) - \sigma(\tau(i))}{\tau(j) -
\tau(i)} \cdot \prod_{i<j} \frac{\tau(j) - \tau(i)}{j - i}
= \varepsilon(\sigma)\,\varepsilon(\tau),
\]
والجداء الأوسط يساوي $\varepsilon(\sigma)$ بعد إعادة الترقيم حسب
الأزواج $\{\tau(i), \tau(j)\}$ (فكل زوج غير مرتب يظهر مرة واحدة،
ويتغير البسط والمقام في الإشارة معًا). وللمبادلة
$\tau = (a\,b)$ مع $a < b$ عدد فردي من الانقلابات؛
وبالعدّ الدقيق: الأزواج المنقلبة $(i, j)$ و$i < j$، مع
$\tau(i) > \tau(j)$، هي
\[
(a, j) \ \text{من أجل } a < j < b, \qquad
(i, b) \ \text{من أجل } a < i < b, \qquad
(a, b) \ \text{نفسه},
\]
أي $(b - a - 1) + (b - a - 1) + 1 = 2(b - a) - 1$ منها،
وهو عدد فردي. (أو بديلًا: نتحقق مباشرة من $(1\,2)$، وفيها انقلاب
واحد، ثم نُقارن --- فللمقارنات إشارة واحدة
لأن $\varepsilon$ تشاكل نحو زمرة تبادلية.) ومنه
$\varepsilon((a\,b)) = (-1)^{2(b-a)-1} = -1$.

\emph{الوحدانية.} تولّد المبادلات $\mathfrak{S}_n$ (فكل
دورة $(a_1\cdots a_k) = (a_1\,a_k)(a_1\,a_{k-1})\cdots(a_1\,a_2)$،
و\cref{thm:b2:structures:cycles} يُتمّ العمل)؛ والتشاكل نحو
$\{\pm1\}$ يتحدد بقيمه على المولّدات.

\emph{النتائج.} تكتب متطابقة الدورات أعلاه \omterm{def:b2:structures:sn}{الدورة} ذات الطول $k$ جداءَ
$k - 1$ \omterm{def:b2:structures:sn}{مبادلة}: فإشارتها $(-1)^{k-1}$. وأما $A_n$: فالتشاكل
$\varepsilon$ شامل (إذ توجد مبادلات من أجل $n \geq 2$)،
و``الصنفان الجانبيان'' $A_n$ و$(1\,2)A_n$ متساويا القوة ويقسمان
$\mathfrak{S}_n$ (بحجة لاغرانج): ومنه $\abs{A_n} =
\frac{n!}{2}$.
\end{proof}

\begin{example}\label{ex:b2:structures:snexample}
$\sigma = \begin{pmatrix} 1&2&3&4&5&6\\ 3&6&5&4&1&2 \end{pmatrix}
= (1\,3\,5)(2\,6)$: الرتبة $\operatorname{lcm}(3,2) = 6$، والإشارة
$(-1)^{2}\cdot(-1)^{1} = -1$. والإشارة أسرع اختبار
للتماثل في الخلط --- وهي محرّك المحدد في
\cref{ch:b2:linalg}.
\end{example}

\begin{example}[ثلاث طرق إلى إشارة واحدة]\label{ex:b2:structures:threeroads}
ليكن $\sigma \in \mathfrak{S}_5$ يرسل $1, 2, 3, 4, 5$ إلى $3, 5, 4,
1, 2$. \emph{عبر الدورات:} $1 \mapsto 3 \mapsto 4 \mapsto 1$ و$2
\mapsto 5 \mapsto 2$، ومنه $\sigma = (1\,3\,4)(2\,5)$ و
$\varepsilon(\sigma) = (-1)^{2}(-1)^{1} = -1$. \emph{عبر
الانقلابات:} في قائمة القيم $3, 5, 4, 1, 2$ تكون الأزواج غير
المرتبة هي $(3,1)$ و$(3,2)$ و$(5,4)$ و$(5,1)$ و$(5,2)$ و$(4,1)$
و$(4,2)$: أي سبعة، ومنه $(-1)^7 = -1$. \emph{عبر
المبادلات:} $\sigma = (1\,4)(1\,3)(2\,5)$، ثلاثة عوامل،
ومنه $(-1)^3 = -1$. ثلاث حسابات، وتماثل واحد: فالوحدانية في
\cref{thm:b2:structures:signature} تضمن ألا يختلف أي نظام مسك دفاتر
عن غيره --- وهذا بالضبط ما يجعل
$\varepsilon$ صالحة بوصفها لا متغيّرًا (انظر مسألة نهاية الأسبوع).
\end{example}

\begin{remark}[إلى أين تمضي الإشارة من هنا]
الإشارة بذرة ثلاثة حصادات لاحقة: فهي تبني المحدد وقاعدة جدائه في
\cref{ch:b2:linalg}؛ وتُشغّل لا متغيّرات التماثل في الأحاجي
التوافقية (وتحلّ مسألة نهاية الأسبوع في هذا الفصل أحجية الخمسة عشر
بها)؛ والزمر المتناوبة $A_n$ التي تعرّفها تصبح مركزية في مجلد
السنة الثالثة، حيث تشرح بساطتها من أجل $n \geq 5$ لماذا لا تُحلّ
المعادلات من الدرجة $5$ بالجذور.
\end{remark}

\section{الحلقات والمثاليات ومجموعات القسمة}

\begin{definition}[المثالي]\label{def:b2:structures:ideal}
لتكن $A$ حلقة تبادلية. \emph{المثالي}\index{مثالي} $I
\subseteq A$ هو زمرة
جزئية جمعية تحقق $a x \in I$ لكل
$a \in A$ و$x \in I$. ونوى تشاكلات الحلقات مثاليات؛ ويكون $I = A$
إذا وفقط إذا $1 \in I$ إذا وفقط إذا احتوى $I$ وحدة. والمثالي
\emph{المولَّد} بالعنصر $x$ هو $xA = \{xa\}$ (ويُسمى مثاليًا \emph{رئيسيًا}).
\end{definition}

\begin{theorem}[{مثاليات $\Z$ ومثاليات $K[X]$}]\label{thm:b2:structures:principal}
كل \omterm{def:b2:structures:ideal}{مثالي} في $\Z$ هو $n\Z$ من أجل $n \in \N$ وحيد؛ وكل \omterm{def:b2:structures:ideal}{مثالي} في
$K[X]$ (حيث $K$ حقل) هو $P\,K[X]$ من أجل $P$ موحَّد وحيد (أو معدوم).
ومن ثم يوجد القاسم المشترك الأكبر في الحلقتين مع علاقات بيزو: $x\Z +
y\Z = \gcd(x,y)\Z$، وكذلك من أجل كثيرات الحدود.
\end{theorem}

\begin{proof}
من أجل $\Z$ كان هذا مبرهنة الزمر الجزئية في مجلد السنة الأولى (\omterm{def:b2:structures:ideal}{فالمثالي}
زمرة جزئية على الخصوص، و$n\Z$ \omterm{def:b2:structures:ideal}{مثالي}). ومن أجل $K[X]$: ليكن
$I \neq \{0\}$ مثاليًا وليكن $P \in I$ غير معدوم من درجة دنيا،
مُوحَّدًا بعد النظم. من أجل $F \in I$، تعطي القسمة الإقليدية $F = PQ + R$
كتابة $R = F - PQ \in I$ مع $\deg R < \deg P$: فتفرض الأصغرية أن $R
= 0$، ومنه $I = P\,K[X]$. أما الوحدانية: فمولّدان موحَّدان يقسم كل
منهما الآخر. وقولا بيزو هما تساوي \omterm{def:b2:structures:ideal}{المثالي} $x\Z +
y\Z$ (أو نظيره
لكثيرات الحدود) مع \omterm{def:b2:structures:ideal}{المثالي} الرئيسي المولَّد بالقاسم المشترك الأكبر ---
وهو نفسه تعريف القاسم المشترك الأكبر المستعمل في السنة الأولى، وقد صار
الآن قولًا عن المثاليات.
\end{proof}

\begin{example}[قاسم مشترك أكبر لكثيري حدود، بطريقتين]\label{ex:b2:structures:polygcd}
لنحسب $\gcd(X^3 - 1,\ X^2 - 1)$ في $\Q[X]$. \emph{بخوارزمية إقليدس:}
\[
X^3 - 1 = X\,(X^2 - 1) + (X - 1),
\qquad
X^2 - 1 = (X + 1)(X - 1) + 0 ,
\]
فالقاسم المشترك الأكبر هو $X - 1$، ويعطي التعويض الرجوعي علاقة
بيزو
\[
X - 1 = 1\cdot(X^3 - 1) - X\cdot(X^2 - 1).
\]
\emph{بالمثاليات:} \omterm{def:b2:structures:ideal}{المثالي} $(X^3 - 1)\Q[X] + (X^2 - 1)\Q[X]$
رئيسي (\cref{thm:b2:structures:principal})؛ وهو يحتوي $X -
1$ (حسب الصيغة المعروضة) ومحتوى في $(X - 1)\Q[X]$ (لأن
المولّدين ينعدمان عند $1$، ومن ثم فهما مضاعفان لكثير الحدود $X - 1$):
فالمولّد الموحَّد هو $X - 1$. والفكرة الختامية: وجهة نظر
المثاليات تحدد القاسم المشترك الأكبر \emph{دون قسمة} --- فالجذور
المشتركة تحدد موقع \omterm{def:b2:structures:ideal}{المثالي}، وإقليدس لا يفعل سوى التصديق عليه.
\end{example}

\begin{definition}[حلقة القسمة $\Z/n\Z$، من جديد]\label{def:b2:structures:quotientring}
من أجل \omterm{def:b2:structures:ideal}{مثالي} $I$ في $A$، تكون العلاقة $x \sim y \iff x - y \in I$ علاقة
تكافؤ متوافقة مع $+$ و$\times$؛ وترث \omterm{def:b2:structures:quotient}{مجموعة القسمة}
$A/I$ بنية حلقة --- هي \emph{حلقة
القسمة}\index{حلقة القسمة} --- تجعل $\pi \colon A \to A/I$
تشاكلًا نواته $I$. ومن أجل $A = \Z$ و$I = n\Z$ نجد
$\Z/n\Z$ المعروفة من مجلد السنة الأولى، ولكن بخاصيتها الشمولية الآن: فكل
تشاكل يُعدم $I$ يتحلل عبر $A/I$.
\end{definition}

\begin{theorem}[مبرهنة البواقي الصينية، بصيغة الحلقات]\label{thm:b2:structures:crt}
إذا كان $\gcd(m, n) = 1$، فإن التطبيق
\[
\Z/mn\Z \longrightarrow \Z/m\Z \times \Z/n\Z,
\qquad
\overline{x} \longmapsto (x \bmod m,\; x \bmod n)
\]
تماثل حلقات.\index{مبرهنة البواقي الصينية} ومن ثم
$\varphi(mn) = \varphi(m)\varphi(n)$ من أجل $m, n$ أوليين فيما بينهما، و
\[
\varphi(n) = n \prod_{p \mid n} \Bigl(1 - \frac 1p\Bigr)
\quad (p \text{ أولي}).
\]
\end{theorem}

\begin{proof}
التطبيق تشاكل حلقات معرَّف جيدًا (فأوجه التوافق مباشرة). أما التباين:
فإن $x \equiv 0$ بترديد $m$ وبترديد $n$ مع
$\gcd(m,n) = 1$ يفرض $mn \mid x$ (بمبرهنة غاوس). وأما الشمول: فللطرفين
$mn$ عنصرًا، فيكفي التباين (تساوي القوى المنتهية) --- أو صراحةً:
انطلاقًا من علاقة بيزو $um +
vn = 1$، فإن صنف
\[
x = b\,um + a\,vn
\]
يُرسَل إلى $(a \bmod m,\ b \bmod n)$، لأن $vn = 1 - um \equiv 1
\pmod m$ يجعل $x \equiv a \pmod m$، وبالتناظر بترديد $n$
--- وهي الوصفة المستعملة عدديًا في
\cref{ex:b2:structures:crtinverse}. وتوافق الوحدات أزواجَ الوحدات (فوحدات حلقة الجداء
هي أزواج الوحدات)، ومنه $\varphi(mn) =
\varphi(m)\varphi(n)$. ومن أجل قوة أولية، $\varphi(p^k) = p^k -
p^{k-1}$ (فغير الوحدات بترديد $p^k$ هي مضاعفات $p$)؛
وتُركّب الضربية صيغة الجداء.
\end{proof}

\begin{example}[قلب التماثل الصيني]\label{ex:b2:structures:crtinverse}
لنأخذ $m = 8$ و$n = 9$. يُصرَّح بمقلوب التماثل بواسطة العنصرين
\emph{المتساويي القوى}: نبحث عن $u \equiv 1 \pmod 8$ بحيث
$u \equiv 0 \pmod 9$ و$v \equiv 0 \pmod 8$، وعن $v \equiv 1 \pmod
9$. من $u = 9k \equiv 1 \pmod 8$: نجد $k \equiv 1$، ومنه $u = 9$؛ ومن
$v = 8k \equiv 1 \pmod 9$: نجد $-k \equiv 1$ و$k \equiv 8$، ومنه $v =
64$. عندئذٍ يكون صنف $x = 9a + 64b$ بترديد $72$ هو الحل الوحيد
للشرطين $x \equiv a \pmod 8$ و$x \equiv b \pmod 9$: فمن أجل $a =
3$ و$b = 5$ نجد $27 + 320 = 347 \equiv 59 \pmod{72}$ ---
وهي بالضبط القيمة الوسطى الموجودة بالتعويض في
\cref{exo:b2:structures:8}. والفكرة الختامية: يحقق $u$ و$v$
العلاقات $u + v \equiv 1$ و$uv \equiv 0$ و$u^2 \equiv u$ و$v^2
\equiv v$ بترديد $72$؛ وهما صورتا $(1, 0)$ و$(0,
1)$،
وكل تفكيك صيني هو في جوهره تفكيك
للواحد $1$ إلى عناصر متساوية القوى ومتعامدة.
\end{example}

\begin{theorem}[أويلر؛ فيرما من جديد]\label{thm:b2:structures:euler}
تشكّل وحدات $\Z/n\Z$ زمرة رتبتها $\varphi(n)$؛ ومنه من أجل
$\gcd(a, n) = 1$:
\[
a^{\varphi(n)} \equiv 1 \pmod n
\qquad (\text{مبرهنة أويلر\index{مبرهنة أويلر}}),
\]
ومبرهنة فيرما الصغرى هي الحالة $n = p$ أولي، وقد صارت على بعد سطر
واحد من لاغرانج.
\end{theorem}

\begin{proof}
الأصناف القابلة للقلب هي بالضبط أصناف الأعداد الصحيحة الأولية مع $n$
(مجلد السنة الأولى): وعددها $\varphi(n)$، وهي تشكّل زمرة بالضرب.
ولاغرانج (\cref{thm:b2:structures:lagrange}):
كل عنصر مرفوع إلى رتبة الزمرة يساوي المحايد.
\end{proof}

\begin{example}[زمرة وحدات بلا مولّد]\label{ex:b2:structures:units15}
للزمرة $(\Z/15\Z)^*$ عدد عناصره $\varphi(15) = \varphi(3)\varphi(5)
= 8$. هل هي دورية؟ لنحسب الرتب مستعملين التماثل
الصيني $(\Z/15\Z)^* \simeq (\Z/3\Z)^* \times (\Z/5\Z)^*$
(فالوحدة بترديد $15$ زوج من الوحدات): رتبتا العاملين هما
$2$ و$4$، ومنه تقسم رتبة كل عنصر
$\operatorname{lcm}(2, 4) = 4 < 8$ --- فلا عنصر يولّدها.
وبالتحديد:
\[
2^4 = 16 \equiv 1, \qquad
4^2 = 16 \equiv 1, \qquad
7^4 \equiv 1, \qquad
11^2 = 121 \equiv 1, \qquad
14^2 \equiv 1 \pmod{15} :
\]
رتب هي $4, 2, 4, 2, 2$ ولا تبلغ $8$ أبدًا. وقارِن ذلك بالحالة
\cref{exo:b2:structures:10}: فالزمرة $(\Z/p\Z)^*$ \emph{دورية} من أجل
$p$ أولي، لأن زمرة الوحدات تقع هناك داخل حقل.
وتبقى مبرهنة أويلر صالحة بالأس $\varphi(15) = 8$،
لكن الأس الشمولي الحقيقي هنا هو $4$ --- فأويلر حدّ
أعلى، وليس دائمًا الحدّ الأمثل.
\end{example}

\begin{definition}[الجبر]\label{def:b2:structures:algebra}
\emph{الجبر}\index{جبر} على $K$ هو فضاء متجهي $K$ على $A$ مزوَّد
ببنية حلقة ضربها ثنائي الخطية على $K$. أمثلة:
$K[X]$ و$\mathcal{M}_n(K)$ و$\mathcal{L}(E)$، وفضاءات الدوال
$\mathcal{F}(X, K)$، و$\C$ بوصفه جبرًا على $\R$. وتشاكلات الجبور
هي تشاكلات الحلقات الخطية؛ و\emph{تطبيق التقييم} $P \mapsto P(u)$
من $K[X]$ نحو $\mathcal{L}(E)$ (أو $\mathcal{M}_n(K)$) هو المثال
المركزي، وهو محرّك \cref{ch:b2:reduction}.
\end{definition}

\begin{example}[تشاكل تقييم ونواته]\label{ex:b2:structures:evaluation}
لنأخذ $A = \begin{pmatrix}0 & 1\\ 0 & 0\end{pmatrix}$ وتطبيق
التقييم $\varepsilon_A \colon \R[X] \to \mathcal{M}_2(\R)$،
$P \mapsto P(A)$. بما أن $A^2 = 0$، فإن
\[
P(A) = P(0)\,I + P'(0)\,A =
\begin{pmatrix} P(0) & P'(0)\\ 0 & P(0)\end{pmatrix},
\]
(إذ لا يبقى من $P$ إلا الحدّان الثابت والخطي). ومنه
$\ker\varepsilon_A = \{P : P(0) = P'(0) = 0\} = X^2\,\R[X]$: وهو
\omterm{def:b2:structures:ideal}{مثالي} رئيسي، تمامًا كما تتنبأ به \cref{thm:b2:structures:principal}، ومولَّد بكثير الحدود $X^2$
الموحَّد ذي الدرجة الدنيا في النواة --- أي \emph{كثير الحدود الأدنى}
للمصفوفة $A$، بطل \cref{ch:b2:reduction}. والصورة هي \omterm{def:b2:structures:algebra}{الجبر}
التبادلي ذو البعد اثنين $\{aI + bA\}$: فتشاكلات التقييم تقلّص
$\R[X]$ غير المنتهي البعد إلى جبور صغيرة قابلة للحساب.
\end{example}

\begin{remark}[آفاق: ثلاثة ألحان ينبغي الإصغاء إليها]\label{rem:b2:structures:perspectives}
ثلاث أفكار بنيوية من هذا الفصل تتكرر في المجلد كله، في توزيع
أوركسترالي أثقل في كل مرة. \emph{التحليل عبر مجموعة قسمة}
(\cref{def:b2:structures:quotient}): فهو يبني $\Z/n\Z$ هنا،
ويعرّف تطبيقات على مجموعات حلول الجمل الخطية في
\cref{ch:b2:linalg}، ويقوم صامتًا تحت كل حجة من نوع
``معرَّف جيدًا على الأصناف''. \emph{اللامتغيّرات}: فالإشارة
تشاكل نحو $\{\pm1\}$ لا تستطيع أي نقلة مشروعة أن تراوغه ---
والمنطق نفسه يعطي قاعدة جداء المحدد
(\cref{ch:b2:linalg})، ولا تغيّر الأثر بالتشابه،
والمقادير المصونة في \cref{ch:b2:diffeq}.
\emph{العدّ في مقابل بنية}: فلاغرانج يعدّ عبر الأصناف
الجانبية، والبعد يُعدّ عبر الأسس
(\cref{ch:b2:linalg})، والتضاعف يُعدّ عبر درجات
كثيرات الحدود (\cref{ch:b2:reduction})؛ وكلما بدا حدّ ما
معجزًا، فثمة تقسيم أو تدريج يقوم بالعدّ.
\end{remark}

\begin{remark}[مزالق شائعة]\label{rem:b2:structures:pitfalls}
أربعة كلاسيكية. (1) يجب التحقق من أن تطبيقًا على مجموعة قسمة
\emph{معرَّف جيدًا}: فالكتابة ``$\overline x \mapsto$ (صيغة في $x$)''
مشروعة فقط إذا كانت الصيغة ثابتة على الأصناف ---
أي توافق \cref{def:b2:structures:quotient}، وليست شكلية.
(2) العلاقة $\operatorname{ord}(ab) =
\operatorname{lcm}(\operatorname{ord}a, \operatorname{ord}b)$
\emph{خاطئة} عمومًا، حتى من أجل عنصرين متبادلين ($a$
و$a^{-1}$)؛ و\cref{exo:b2:structures:4} يعطي القول الصحيح
في حالة التبادل مع رتبتين أوليتين فيما بينهما، والدورات المنفصلة
تعطي الصيغة الصحيحة للتبديلات. (3) تصمد قابلية العد أمام
\emph{الاتحادات} \omterm{def:b2:structures:countable}{القابلة للعد} و\emph{الجداءات} المنتهية، لكن ليس
أمام الجداءات \omterm{def:b2:structures:countable}{القابلة للعد}: فالمجموعة $\{0,1\}^{\N}$ غير \omterm{def:b2:structures:countable}{قابلة للعد}
(\cref{exo:b2:structures:3}) رغم أن لكل عامل عنصرين
اثنين. (4) لا تحتاج مبرهنة كانتور--برنشتاين إلا إلى غرستين في
الاتجاهين، لكن التقابل الذي تبنيه غير متصل عادة وغير صريح ---
فلا تنتظر صيغة
(\cref{ex:b2:structures:cbexample}).
\end{remark}

\begin{remark}[أين يُستعمَل هذا الفصل]
في كل مكان تقريبًا. فالإشارة تبني المحددات
(\cref{ch:b2:linalg})؛ وتشاكل التقييم $P \mapsto P(u)$
والمثاليات الرئيسية في $K[X]$ تنتج كثيرات الحدود الدنيا
وتفكيكات النواة في \cref{ch:b2:reduction}؛ وقابلية العد
هي المسرح الذي يؤدي عليه \cref{ch:b2:proba} (الاحتمال على الفضاءات
\omterm{def:b2:structures:countable}{القابلة للعد}) والسبب في أن الطوبولوجيا لا تكفّ عن إنتاج مجموعات
كثيفة \omterm{def:b2:structures:countable}{قابلة للعد} (\cref{ch:b2:metric}). ويُعاد نشر إنشاء القسمة $A/I$
في مجلد السنة الثالثة لبناء الحقول $K[X]/(P)$، ومنها نظرية غالوا:
فالخاصية الشمولية المبرهَنة هنا تُستعمل هناك حرفيًا.
\end{remark}

\section{تمارين}

\begin{exercise}[$\star$]\label{exo:b2:structures:1}
أي المجموعات التالية \omterm{def:b2:structures:countable}{قابلة للعد}؟ مجموعة الأجزاء المنتهية
من $\N$؛ ومجموعة \emph{كل} أجزاء $\N$؛ و$\R \setminus \Q$؛
ومجموعة كثيرات الحدود ذات المعاملات الناطقة؛ ومجموعة
المتتاليات المكوَّنة من $0$ و$1$ والمنعدمة ابتداءً من رتبة ما.
\end{exercise}

\begin{exercise}[$\star$]\label{exo:b2:structures:2}
في $\mathfrak{S}_7$، ليكن $\sigma = (1\,4\,2\,6)(3\,5)$ و$\tau =
(2\,3\,7)$. احسب $\sigma\tau$ و$\tau\sigma$ على صورة دورات
منفصلة، ورتب التبديلات الأربع وإشاراتها، و$\sigma^{2026}$.
\end{exercise}

\begin{exercise}[$\star$]\label{exo:b2:structures:3}
أنشئ غرسات صريحة تُبيّن أن $\mathcal{P}(\N)$ و$\intcc{0}{1}$ ومجموعة المتتاليات
الثنائية $\{0,1\}^{\N}$ متساوية القوة مثنى مثنى \emph{(بالنشر الثنائي في
الاتجاهين؛ وتستوعب مبرهنة كانتور--برنشتاين إزعاج التمثيل المزدوج)}.
\end{exercise}

\begin{exercise}[$\star$]\label{exo:b2:structures:4}
لتكن $G$ زمرة وليكن $a, b \in G$ عنصرين متبادلين رتبتاهما المنتهيتان
$m$ و$n$ أوليتان فيما بينهما. برهن على أن $\operatorname{ord}(ab) =
mn$. بيّن بمثال في $\mathfrak{S}_3$ أن التبادل شرط جوهري.
\end{exercise}

\begin{exercise}[$\star\star$]\label{exo:b2:structures:5}
لتكن $G$ زمرة منتهية رتبتها زوجية. برهن على أن $G$ تحتوي
عنصرًا رتبته $2$. \emph{(قابِل بين كل عنصر ومقلوبه؛
وعُدّ العناصر المقابَلة بنفسها.)}
\end{exercise}

\begin{exercise}[$\star\star$]\label{exo:b2:structures:6}
برهن على أن $A_n$ (من أجل $n \geq 3$) مولَّدة بالدورات ذات الطول $3$.
\emph{(جداء مبادلتين هو دورة ذات الطول $3$ أو جداء دورتين ذواتَي
الطول $3$.)}
\end{exercise}

\begin{exercise}[$\star\star$]\label{exo:b2:structures:7}
عيّن كل تشاكلات الزمر: من $(\Q, +)$ نحو $(\Z, +)$؛ ومن
$(\Z/n\Z, +)$ نحو $(\Z/m\Z, +)$ \emph{(عُدّها:
$\gcd(m,n)$)}؛ ومن $(\Q, +)$ نحو $(\Q_+^*, \times)$.
\end{exercise}

\begin{exercise}[$\star\star$]\label{exo:b2:structures:8}
باستعمال مبرهنة البواقي الصينية، احسب $\varphi(360)$، وجد
كل $x$ بحيث $x \equiv 3 \pmod 8$ و$x \equiv 5 \pmod 9$ و$x
\equiv 2 \pmod 5$، واحسب الرقمين الأخيرين من $3^{2026}$
\emph{(أويلر بترديد $100$؛ وانتبه: اعمل بترديد $4$ وبترديد $25$)}.
\end{exercise}

\begin{exercise}[$\star\star\star$]\label{exo:b2:structures:9}
برهن على أن كل حلقة تامة منتهية حقل. استنتج أن
$\Z/n\Z$ حقل إذا وفقط إذا كان $n$ أوليًا (مرة أخرى).
\end{exercise}

\begin{exercise}[$\star\star\star$]\label{exo:b2:structures:10}
(كلاسيكي) ليكن $K$ حقلًا ولتكن $G$ زمرة جزئية \emph{منتهية} من
$(K^*, \times)$. برهن على أن $G$ دورية.
\emph{إرشاد: ليكن $m$ أكبر رتبة بين عناصر $G$؛
بيّن أن رتبة كل عنصر تقسم $m$ (باستعمال
\cref{exo:b2:structures:4} على أجزاء أولية فيما بينها مناسبة)، فتحقق $G$ كلها
$x^m = 1$؛ ثم عُدّ جذور $X^m - 1$.}
وعلى الخصوص فإن $(\Z/p\Z)^*$ دورية.
\end{exercise}

\begin{exercise}[$\star\star\star$]\label{exo:b2:structures:11}
برهن على أن الزمرة $(\Q, +)$ ليست دورية، والأسوأ: أنها ليست
حتى مولَّدة بجزء منته. وبرهن في المقابل على أن كل زمرة جزئية من
$(\Q, +)$ مولَّدة بجزء منته تكون دورية.
\end{exercise}

\begin{exercise}[$\star\star$]\label{exo:b2:structures:12}
(محك ديدكيند) برهن على أن كل مجموعة غير منتهية تحتوي جزءًا
قابلًا للعد، واستنتج أن مجموعة $E$ غير منتهية إذا وفقط
إذا كانت مساوية القوة لجزء صحيح منها. \emph{(من أجل الاتجاه
المباشر، أزِح جزءًا قابلًا للعد خطوة واحدة؛ ومن أجل العكس،
تذكّر مبدأ الجحور.)}
\end{exercise}

\section{مسألة: أحجية الخمسة عشر}

أحجية الخمسة عشر لوحة من الحجم $4 \times 4$ تحمل خمس عشرة قطعة منزلقة
مرقّمة من $1$ إلى $15$ وخانة فارغة واحدة؛ وتنقل النقلةُ إحدى
القطع المجاورة للخانة الفارغة إليها. وفي تسعينيات القرن التاسع عشر
روّج سام لويد للأحجية بعرضه \$1000 لمن يستطيع أن
يبادل بين القطعتين $14$ و$15$ ويعيد كل قطعة أخرى إلى
مكانها. ولم يقبض أحد الجائزة قط، وتبرهن مسألة نهاية الأسبوع هذه على
شقّي السبب: فإشارة
\cref{thm:b2:structures:signature} تمنع \omterm{def:b2:structures:sn}{مبادلة} لويد،
و--- وهذا هو الشق الأصعب، الإنشائي --- \emph{كل} ما تسمح به الإشارة
قابل للحلّ فعلًا. والصياغة الكاملة هي مبرهنة جونسون--ستوري (1879).

\begin{omfigure}
\begin{tikzpicture}[scale=0.72]
  \draw[omDef, thick] (0,0) grid (4,4);
  \draw[omDef, very thick] (0,0) rectangle (4,4);
  \foreach \k in {1,...,15} {
    \pgfmathtruncatemacro\row{int((\k-1)/4)}
    \pgfmathtruncatemacro\col{int(mod(\k-1,4))}
    \node at (\col+0.5, 3.5-\row) {$\k$};
  }
  \node[below, font=\small] at (2,-0.2) {الوضعية المحلولة};
\end{tikzpicture}
\qquad
\begin{tikzpicture}[scale=0.72]
  \draw[omDef, thick] (0,0) grid (4,4);
  \draw[omDef, very thick] (0,0) rectangle (4,4);
  \foreach \k in {1,...,13} {
    \pgfmathtruncatemacro\row{int((\k-1)/4)}
    \pgfmathtruncatemacro\col{int(mod(\k-1,4))}
    \node at (\col+0.5, 3.5-\row) {$\k$};
  }
  \node[omThm] at (1.5,0.5) {$15$};
  \node[omThm] at (2.5,0.5) {$14$};
  \node[below, font=\small] at (2,-0.2) {جائزة لويد};
\end{tikzpicture}

{\small الوضعية المحلولة ووضعية $14$--$15$ عند سام لويد.
والسؤال بجائزة \$1000: هل تستطيع نقلات مشروعة أن تحوّل اللوحة اليمنى
إلى اليسرى؟}
\end{omfigure}

\begin{problem}[{مسألة نهاية الأسبوع --- مبرهنة جونسون--ستوري
في قابلية الحلّ}]
\label{pb:b2:structures:1}
رقّم الخانات من $1$ إلى $16$ بترتيب القراءة (من اليسار إلى اليمين،
ومن الأعلى إلى الأسفل)، بحيث تقع الخانة $k$ في السطر $i$ والعمود $j$
مع $k
= 4(i - 1) + j$. والخانة $16$ (أسفل اليمين) هي \emph{موطن}
الخانة الفارغة؛ ونعامل الخانة الفارغة كقطعة سادسة عشرة،
تُكتب $b$ وتُماهى مع العدد $16$. وتكون
\emph{الوضعية} تقابلًا $\sigma \colon \intint1{16}
\to \intint1{16}$، أي خانة $\mapsto$ محتواها؛ والوضعية
\emph{المحلولة} هي $\sigma = \mathrm{id}$. وفي كل ما يلي،
$\varepsilon$ هي إشارة
\cref{thm:b2:structures:signature}، وتكون خانتان
\emph{متجاورتين} إذا اشتركتا في حافة من اللوحة.

\medskip
\noindent\textbf{الجزء الأول --- الوضعيات والنقلات والإشارات.}
\begin{enumerate}
  \item برّر أن الوضعيات هي بالضبط عناصر
        $\mathfrak{S}_{16}$، فيكون عددها $16! =
        20\,922\,789\,888\,000$، وأن عدد
        النقلات المشروعة انطلاقًا من وضعية معطاة هو $2$ أو $3$ أو $4$،
        تبعًا لوقوع الخانة الفارغة في زاوية أو على حافة أو
        في الداخل.
  \item لتكن $\sigma$ وضعية، ولتكن $p = \sigma^{-1}(16)$
        خانة الفراغ، ولتكن $c$ خانة مجاورة للخانة $p$. بيّن
        أن انزلاق قطعة $c$ إلى $p$ يعطي
        الوضعية $\sigma' = \sigma \circ \tau$ مع $\tau =
        (p\ c)$، واستنتج أن كل نقلة تقلب الإشارة:
        $\varepsilon(\sigma') = -\varepsilon(\sigma)$.
  \item لوّن اللوحة كرقعة شطرنج: $\chi(k) = (-1)^{i+j}$ من أجل الخانة
        $k$ في السطر $i$ والعمود $j$. بيّن أن كل نقلة تقلب
        $\chi(\text{خانة الفراغ})$، واستنتج أن كل
        متتالية نقلات تعيد الفراغ إلى خانته الابتدائية
        يكون طولها زوجيًا.
  \item بيّن أن
        \[
        I(\sigma) = \varepsilon(\sigma)\,
        \chi\bigl(\sigma^{-1}(16)\bigr)
        \]
        لا متغيّر تحت كل نقلة مشروعة، واحسب
        $I(\mathrm{id})$.
\end{enumerate}

\medskip
\noindent\textbf{الجزء الثاني --- جائزة لويد: اللامتغيّر في
العمل.}
\begin{enumerate}[resume]
  \item توافق وضعية لويد $\sigma_L$ الوضعية المحلولة
        إلا في أن الخانتين $14$ و$15$ تحملان القطعتين $15$ و$14$.
        احسب $I(\sigma_L)$ واستنتج أن لا متتالية نقلات
        تربط $\sigma_L$ بالوضعية المحلولة:
        فجائزة لويد لم تكن في خطر قط.
  \item بيّن أن نصف الوضعيات بالضبط يحقق $I =
        +1$: أي $\abs{\{\sigma : I(\sigma) = +1\}} = 16!/2$.
        \emph{(من أجل خانة فراغ ثابتة، قابِل بين الوضعيات
        بالتركيب مع \omterm{def:b2:structures:sn}{مبادلة} ثابتة لخانتين أخريين.)}
  \item بيّن أن كل نقلة تُلغى بنقلة مشروعة، وأن ``$\sigma'$ قابلة
        للبلوغ من $\sigma$ بنقلات مشروعة'' علاقة تكافؤ، وأن
        صنف $R$ الوضعية المحلولة يحقق $R \subseteq \{I = +1\}$.
        استنتج أن هناك صنفين على الأقل.
  \item لنفترض أن الفراغ في موطنه: $\sigma(16) = 16$. بيّن أن
        $I(\sigma) = \varepsilon(\rho)$ حيث $\rho \in
        \mathfrak{S}_{15}$ مقصور $\sigma$ على
        الخانات $1, \dots, 15$، وأن كل وضعية يمكن
        نقلها بنقلات مشروعة إلى وضعية فراغها في موطنه.
        استنتج: لبرهان $R = \{I = +1\}$ يكفي أن نحقق
        كل تبديلة \emph{زوجية} للخانات الخمس عشرة غير الموطن
        بمتتالية نقلات تبدأ وتنتهي والفراغ في موطنه.
\end{enumerate}

\medskip
\noindent\textbf{الجزء الثالث --- جولات الفراغ وزمرة
البرامج.} \emph{البرنامج} متتالية منتهية من النقلات المشروعة،
تبدأ من وضعية فراغها في موطنه، وتكون وضعيتها النهائية
بفراغ في موطنه أيضًا. و\emph{أثره} هو
التبديلة $\pi$ للخانات المعرَّفة بما يلي: محتوى الخانة $x$
ينتهي في الخانة $\pi(x)$.
\begin{enumerate}[resume]
  \item بيّن أن برنامجًا يُنفَّذ انطلاقًا من $\sigma$ ينتهي عند $\sigma
        \circ \pi^{-1}$؛ وأن تنفيذ برنامجين متتاليين
        يركّب أثريهما؛ وأن مجموعة $H$ كل الآثار زمرة
        جزئية من $\mathfrak{S}_{15}$ (تبديلات
        الخانات $1, \dots, 15$) محتواة في الزمرة
        المتناوبة $A_{15}$.
  \item (الجولة الأولية) انطلاقًا من الفراغ في موطنه، أزلق
        الفراغ حول المربّع $2 \times 2$ أسفل اليمين: أي الخانات
        $16 \to 12 \to 11 \to 15 \to 16$. بيّن أن الأثر هو
        \omterm{def:b2:structures:sn}{الدورة} ذات الطول $3$، أي $(11\ 12\ 15)$، وأن الجولة العكسية تعطي
        $(11\ 15\ 12)$. وكلاهما في $H$.
  \item (الجولة الكبرى) تحقّق من أن
        \[
        16 \to 15 \to 14 \to 13 \to 9 \to 5 \to 1 \to 2 \to 3
        \to 4 \to 8 \to 7 \to 6 \to 10 \to 11 \to 12 \to 16
        \]
        مسير مغلق يمرّ بالخانات الست عشرة كلها (بخطوات بين
        متجاورين فقط)، وأن أثره هو \omterm{def:b2:structures:sn}{الدورة} ذات الطول $15$
        \[
        \zeta = (15\ 12\ 11\ 10\ 6\ 7\ 8\ 4\ 3\ 2\ 1\ 5\ 9\ 13\
        14) .
        \]
        وبكتابة $x_0 = 15$ و$x_1 = 12$ و$x_2 = 11$، \dots، و$x_{14}
        = 14$ لترتيبها الدوري، تحقّق من أن الجولة
        الأولية العكسية في السؤال 10 هي بالضبط $(x_0\ x_1\
        x_2)$.
  \item برهن على صيغة المقارنة في أي $\mathfrak{S}_n$: من أجل
        تبديلة $g$ ودورة ذات الطول $3$،
        \[
        g\,(a\ b\ c)\,g^{-1} = \bigl(g(a)\ g(b)\ g(c)\bigr),
        \]
        ولاحظ أن $H$، بوصفها زمرة، مغلقة تحت
        المقارنة بعناصرها.
  \item استنتج أن $H$ تحتوي كل الدورات \emph{المتتالية}
        الخمس عشرة ذات الطول $3$ في الجولة الكبرى:
        \[
        s_t = (x_t\ x_{t+1}\ x_{t+2}) \qquad (t \in \Z/15\Z,
        \text{ بأدلة بترديد } 15).
        \]
\end{enumerate}

\medskip
\noindent\textbf{الجزء الرابع --- توليد الزمرة المتناوبة.}
\begin{enumerate}[resume]
  \item (المبرهنة المساعدة أ) لتكن $s$ و$t$ دورتين ذواتَي الطول $3$
        يشترك حاملاهما في نقطتين بالضبط، وليكن حاملاهما $\{a, b, c\}$ و
        $\{b, c, d\}$. بيّن أنه بعد تعويض $s$ أو $t$
        بمقلوبها عند الاقتضاء (وهذا لا يغيّر شيئًا في الزمرة
        الجزئية \omterm{def:b2:structures:generated}{المولَّدة})، يكون الجداء $st$ \omterm{def:b2:structures:sn}{مبادلة}
        مضاعفة؛ وبيّن أن $A_4$ لا تحتوي زمرة جزئية
        رتبتها $6$ \emph{(فالزمرة الجزئية ذات الدليل $2$ تحتوي كل
        مربّع؛ عُدّ الدورات ذات الطول $3$ بين المربعات)}؛ واستنتج
        أن $\langle s, t\rangle$ هي الزمرة المتناوبة كاملةً
        على الحروف الأربعة $\{a, b, c, d\}$.
  \item (المبرهنة المساعدة ب) لتكن $X$ مجموعة من $k \geq 4$ حرفًا، مع $w
        \notin X$،
        ولتكن $G$ زمرة جزئية من $\mathfrak{S}_n$ ما تحتوي كل تبديلة
        زوجية للمجموعة $X$ ودورة واحدة ذات الطول $3$ هي $(u\ v\ w)$ مع $u, v \in X$. بيّن
        أنه من أجل كل $a, b \in X$ مختلفة توجد تبديلة \emph{زوجية}
        $g$ للمجموعة $X$ تحقق $g(u) = a$ و$g(v) = b$، ثم
        استنتج أن $(a\ b\ w) \in G$.
  \item استنتج أن الزمرة $G$ في المبرهنة المساعدة ب تحتوي كل تبديلة
        زوجية للمجموعة $X \cup \{w\}$ \emph{(استعمل
        \cref{exo:b2:structures:6}: فالدورات ذات الطول $3$ تولّد)}.
        ثم، بتسلسل المبرهنتين المساعدتين أ وب على الدورات
        المتتالية ذوات الطول $3$، أي $s_0, s_1, \dots, s_{12}$، في السؤال 13،
        برهن على أن $\langle s_0, \dots, s_{12}\rangle = A_{15}$.
  \item استنتج أن $H = A_{15}$: \emph{أي أن كل إعادة ترتيب زوجية
        للقطع الخمس عشرة قابلة للتحقيق ببرنامج}، وأن $H$
        تحتوي $15!/2 = 653\,837\,184\,000$ عنصرًا.
  \item (مبرهنة جونسون--ستوري، 1879) اجمع الأسئلة 6 و7
        و8 و17: الوضعيات القابلة للبلوغ من الوضعية المحلولة
        هي \emph{بالضبط} الوضعيات $16!/2 =
        10\,461\,394\,944\,000$ التي تحقق $I = +1$؛
        وللبلوغ صنفان \emph{اثنان} بالضبط، صنف
        الوضعية المحلولة وصنف وضعية لويد
        $\sigma_L$. \emph{(من أجل النقطة الثانية، أعِد تسمية القطعتين
        $14$ و$15$: بيّن أن $\sigma \mapsto (14\ 15) \circ
        \sigma$ يرسل متتاليات النقلات إلى متتاليات نقلات
        ويبادل بين $\{I = +1\}$ و$\{I = -1\}$.)}
\end{enumerate}

\medskip
\noindent\textbf{الجزء الخامس --- المحكّات والمتغيّرات
والنظرة من فوق.}
\begin{enumerate}[resume]
  \item (المحك العملي) اقرأ القطع الخمس عشرة بترتيب قراءة
        خاناتها، متجاوزًا الفراغ، وليكن $N$ عدد انقلابات هذه
        القائمة؛ وليكن $r$ سطر الفراغ محسوبًا من \emph{الأسفل}. بيّن
        أن $I(\sigma) = (-1)^{N + r + 1}$، فتكون $\sigma$
        قابلة للحلّ إذا وفقط إذا كان $N + r$ فرديًا.
  \item (فعل الزمر) \emph{فعل} زمرة $G$ على مجموعة
        $X$ هو تطبيق $G \times X \to X$، $(g, x) \mapsto g \cdot
        x$، يحقق $e \cdot x = x$ و$g \cdot (h \cdot x) =
        (gh) \cdot x$؛ و\emph{مدار} $x$ هو $G \cdot x$، ويكون
        الفعل \emph{حرًا} إذا كان $g \cdot x = x$ يفرض $g =
        e$. بيّن أن $h \cdot \sigma = \sigma \circ h^{-1}$
        يعرّف فعلًا حرًا للزمرة $H$ على مجموعة الوضعيات ذوات
        الفراغ في الموطن، وأن مداراته هي بالضبط أصناف
        البلوغ المتبادل بالبرامج، واستنتج من عدد المدارات
        أن هذه الوضعيات تنقسم إلى $15!\,/\,\abs H = 2$ صنفًا بالضبط.
  \item (عائق $3 \times 3$) بيّن أن اللوحة $3 \times 3$
        لا تقبل \emph{أي} مسير مغلق يزور كل خانة
        مرة واحدة بالضبط: فتفشل استراتيجية الجولة الكبرى
        في الجزء الثالث من أجل أحجية الثمانية. \emph{(لوّن الخانات
        التسع كرقعة شطرنج.)}
  \item (الإصلاح) على اللوحة $3 \times 3$ ذات الخانات من $1$ إلى
        $9$ بترتيب القراءة وموطنٍ هو $9$: احسب أثري
        جولة المحيط $9 \to 8 \to 7 \to 4 \to 1 \to 2 \to 3
        \to 6 \to 9$ (وهي دورة ذات الطول $7$، أي $\zeta'$، تثبّت المركز $5$)
        وجولة الزوايا $9 \to 6 \to 5 \to 8 \to 9$ (وهي
        دورة ذات الطول $3$ تمرّ بالمركز). وبمقارنة الثانية بقوى
        $\zeta'$ وتسلسل المبرهنتين المساعدتين أ وب، برهن
        على أن زمرة برامج أحجية الثمانية هي $A_8$ كاملةً،
        ومن ثم على أن $9!/2 = 181\,440$ بالضبط من الوضعيات $9! =
        362\,880$ قابلة للحلّ.
  \item (لوحة فقيرة) لتكن اللوحة الآن دورة واحدة من $n
        \geq 4$ خانة
        تحمل $n - 1$ قطعة. بيّن أن الترتيب الدوري
        للقطع لا متغيّر، وأن كل صنف بلوغ يحتوي بالضبط $n(n - 1)$
        وضعية \emph{(فالأصناف هي مدارات \omterm{def:b2:structures:generated}{زمرة دورية} رتبتها
        $\operatorname{lcm}(n, n-1) = n(n-1)$)}، وأن عدد الأصناف
        هو $(n - 2)!$ --- وهو من أجل $n \geq 5$ أكبر بكثير من
        $2$: فعلى لوحة نحيلة لا يلتقط لا متغيّر التماثل شيئًا
        يُذكر، والهندسة هي الحاكمة.
  \item حكمان بمحك السؤال 19: اللوحة المقلوبة كليًا
        (القطع $15, 14, \dots, 1$ في الخانات $1$
        إلى $15$، والفراغ في موطنه)، واللوحة التي فيها الفراغ في
        الخانة $1$ تتلوه القطع $15, 14, \dots, 1$ في الخانات $2$
        إلى $16$. أيهما قابلة للحلّ؟
  \item (تركيب) للبرهان ركيزتان مستقلتان: \emph{لا متغيّر}
        ($I$، مبني على تشاكل الإشارة)
        يبيّن أن نصف الوضعيات على الأكثر قابل للبلوغ،
        ومبرهنة \emph{توليد صريح} ($H = A_{15}$)
        تبيّن أن نصفها على الأقل قابل للبلوغ. بجملة واحدة لكل
        منها، قل أين دخل ما يلي: خاصية التشاكل في
        $\varepsilon$؛ ومبرهنة لاغرانج؛ وتوليد $A_n$
        بالدورات ذات الطول $3$؛ والمقارنة. وصُغ المبدأ الشامل في
        سطر واحد.
\end{enumerate}
\end{problem}
